\section{The ecology, the frozen model, and the individuals}
\label{sec:methods}

We study whether a frozen pretrained decision model, used as a source of prior judgements about food descriptions, helps individuals whose decision-layer parameters are evolved and whose beliefs are learned within a lifetime. The design goal is that three sources of competence can be separated: \emph{innate} parameters (the genome of an individual), \emph{acquired} knowledge (beliefs formed from the individual's own experience), and \emph{perception} supplied by a frozen model whose weights are never changed. Evolution of the genome is not evolution of the frozen model.

\subsection{The ecology}
\label{sec:ecology}
A shared simulator owns the world state. Individuals (the cohort) only submit action proposals; energy, position and survival are changed exclusively by the simulator, and an energy ledger (initial energy plus supply equals living energy, stored food, and every recorded loss) is verified after each step in the test suite. World time does not depend on rendering. Randomness comes from simulator-managed generators, and a run is reproducible from its configuration and the recorded actions.

Unless stated otherwise, the world is a $16\times 12$ grid with eight individuals and a horizon of 3600 ticks, with one decision every ten ticks (360 rounds). An individual has capacity 40, starts with 32 units of energy and pays a resting metabolism of 0.012 per tick and a movement cost per unit distance. Food (140 units initially, 14 units added every ten ticks to random cells) decays at 0.003 per tick; there is no diffusion, no reproduction and no return of energy at death, so the cohort is fixed and the energy budget is closed. An individual sees only the cells within a perception radius of three, a memory of recently seen food (at most eight entries), and its own state; it never sees global positions or hidden properties of other individuals. A small discrete action set (move, eat, rest) is available. Eating transfers at most 0.6 units per tick from the current cell.

\paragraph{Semantic food.}
Each food-bearing cell carries one \emph{kind}. The kind has a natural-language description that is visible to individuals and a hidden effect factor $f$: eating an amount $a$ changes the individual's energy by $af$. Factors are drawn from four tiers: nutritive ($f=+1$), mild ($-0.4$), strong ($-1$) and lethal ($-4$); in all experiments reported here the lethal tier is allowed (damage cap 4). A world samples $K$ kinds (32 by default), half nutritive and half toxic, where, among toxic kinds, 15\% are lethal, 40\% strong and 45\% mild. \emph{Deception} flips the true effect of a fraction of kinds ($25\%$ or $40\%$) without changing their description: a nutritive description may hide a harmful effect, and a toxic description a nutritive one. Hidden factors are used only by the simulator to settle energy and by the oracle (Section~\ref{sec:individuals}) to measure ceilings; they never enter any individual's input, feature or training signal.

\paragraph{Descriptions and dialects.}
Descriptions are produced by a generator that is independent of any model output. Each kind has a hidden tier, and its text is composed of two or three clauses drawn from five fields (appearance, smell, feel, place, signs). A clause is a positive, negative, or neutral phrase; with probability 0.9 it agrees with the truth, and 20\% of clauses are replaced by an uninformative phrase. We wrote four dialects with disjoint phrasing: a \emph{training} dialect (12,000 descriptions), a \emph{test} dialect (4,000), a \emph{held-out} field-notes dialect (4,000, written blind to any model output) and a \emph{factual} dialect (4,000, neutral in tone, whose clues are facts and indirect knowledge such as related species, habitat, animal behaviour and the effect of cooking). For the factual dialect we also generated a separate 12,000-description training pool that shares the clause bank with, but contains no text identical to, the evaluation pool. All dialects were written by one author; this is a limitation (Section~\ref{sec:limitations}).

\paragraph{Delayed consequences.}
With a delay of $d$ ticks, the energy effect of a bite lands only after $d$ ticks. Until then the gain is recorded as \emph{in transit} in the ledger (and is lost if the individual dies first), and when the effect lands the individual perceives a one-sentence sensation (for example, ``a dull ache stirs in your belly''). We use $d\in\{0,20,40,80\}$.

\paragraph{Ripening and seasons (language cues that change over time).}
In the ripening world the effect of a nutritive kind depends on a stage that cycles with time: the season is $\lfloor t/300\rfloor \bmod 4$ and each kind has a fixed phase offset. The four stages multiply the factor of nutritive kinds by $0.3$ (unripe), $1$ (ripe), $0.4$ (overripe) and $-0.8$ (rotten); toxic kinds do not change. The description shown to the individual is the base description followed by a stage phrase (for example, ``gone soft and sour''), and the observation carries a season sentence that never states a value. Two phrasings of the four stage phrases exist: set A is used for training and set B for evaluation. A cue can be \emph{true}, \emph{absent} (the value still follows the stage) or \emph{mismatched} (the phrase of stage $s+2$ is shown), which changes observations but never the settlement. We use $K=16$ (a hard world) and $K=4$ (a world that experience-only learners can solve; Section~\ref{sec:results}).

\paragraph{Social observation.}
When enabled, a visible neighbour is accompanied by a one-sentence description of its appearance (for example, ``doubled over, clutching its belly''), chosen from a bank according to its most recent noticeable consequence (at least 0.3 energy, about half a bite; when delayed, at the time the effect lands), and by the descriptions of up to three foods it ate in the last 300 ticks. Appearances persist for 60 ticks. Individuals do not share parameters, beliefs or internal experience. Appearance phrases also come in two sets (A for training, B for evaluation) and can be true, absent or mismatched (benefit and harm swapped).

\subsection{The frozen model as ``System 1''}
\label{sec:laya}
The frozen model is \laya{} (a ModernBERT-large encoder with a fine-tuned two-layer decision head, about 421M parameters) \unv{Laya model card and the Jev/Laya documentation}. It is used only as a source of cached judgements and representations: for each food description we ask whether the creature should eat the food, in two option orders, and average the probability of ``eat'' (the head has a strong preference for the first option, which the two-order average removes). The result, a probability $p$ per description text, is cached and read by individuals as a prior; the model is not called in the simulation loop. We also cache the encoder's mean-pooled representation of each description, projected to 32 dimensions with a PCA fitted on the training dialect only, for the experiments that use embeddings. All cached values depend on the text alone.

\subsection{Individuals and genomes}
\label{sec:individuals}
Individuals are \emph{structured}: a fixed rule skeleton moves toward allowed food and eats (a rule that was validated before any evolution), and the genome controls how the individual treats priors and experience. An individual keeps a belief for each description it has eaten, $(n,s)$ (the number of experiences and their sum), updated with the realised energy gain $x=\mathrm{clip}(\mathrm{gain}/3,-1,1)$ and with a forgetting rate. The value of a description $d$ is
\begin{equation}
v(d)=\frac{n_0\,v_0(d)+s}{n_0+n}+\frac{c}{1+n},\qquad v_0(d)=\mathrm{clip}\bigl(a\,(2p_d-1)+b_0\bigr),
\end{equation}
where $n_0=e^{g_1}$ is the trust in the prior (a pseudo-count), $c$ a curiosity bonus, and $a,b_0$ a calibration slope and bias applied to the prior $p_d$. Foods with $v<\tau$ are vetoed (hidden from the skeleton), except for a retry probability. Further genes set the forgetting rate, the minimum energy fraction required before trying an unseen food, and, in delayed worlds, a \emph{credit window}: when a large energy shock arrives, the individual credits it to every kind eaten in the preceding window (up to 150 ticks). The base genome has nine genes; social individuals have one more (the weight of social evidence, Section~\ref{sec:ecology}).

\paragraph{Prior conditions.}
The \emph{prior source} defines $p_d$: \texttt{laya} (the cached judgement above), \texttt{none} ($v_0=0$, so the individual learns from experience only), \texttt{vader} (the valence lexicon VADER \unv{Hutto and Gilbert, VADER}, squashed to $[0,1]$, as an independent sentiment baseline), and \texttt{shuffled} (the cached judgements assigned to other descriptions, which preserves their distribution but destroys the correspondence). A hand-written oracle that reads the hidden factors, and a blind rule that eats everything it reaches, are used only to bound the achievable range; the oracle never produces an individual.

\paragraph{Social individuals.}
A social individual wraps each neighbour appearance as a sentence about a food (``a food that a nearby creature just ate; the creature now looks \dots'') and reads it with the same prior source as the food descriptions. When the reading is clearly favourable or unfavourable, the neighbour's recent meals receive one weighted piece of evidence (the weight is the social gene); the same appearance and meals count once within 120 ticks, because neighbours have no identity across decisions. An individual with no prior cannot read the sentence and ignores neighbours.

\subsection{Evolution}
\label{sec:evolution}
Genomes are optimised with an evolution strategy \unv{Salimans et al., OpenAI-ES}: antithetic sampling, rank-based fitness shaping, an Adam update, normalised gene coordinates, population 24, perturbation scale 0.25 and learning rate 0.1, for 40 generations. Each generation evaluates every candidate on six training worlds drawn from the training dialect (the deception rate alternates between 0.25 and 0.40 where stated), and the cohort of eight identical individuals shares the genome. Fitness is $0.7\times$(fraction of agent-ticks survived)$+0.3\times$(final energy as a fraction of capacity), averaged over the cohort. The training world sequence is the same for every evolutionary run; runs differ in the noise of the strategy, and the ``training seed'' pairs the same noise across conditions. \emph{Matched} training means that training and evaluation share the consequence delay (and, where relevant, the cue structure).
